\manualchapter{Control reference}{sec:controls}
The upper and lower slider rows belong to lanes 1 and 2. Within each row,
read Amplitude, Attack, Decay, Sustain Level and Sustain Rate left to right.
The bottom socket groups provide each lane's Gate, Retrig, Envelope and
Inverted connections. Timing sliders have no CV inputs.

\begin{tabularx}{\textwidth}{@{}l l X@{}}
\toprule Slider & Range; default & Implemented behavior \\\midrule
Amplitude & $-128$--127; 127 & Signed 8-bit gain; zero silences, negative values invert. \\
Attack & 0--15; 10 & Reversed before reaching the chip; larger settings are slower. \\
Decay & 0--7; 7 & Reversed before reaching the chip; larger settings are slower. \\
Sustain Level & 0--7; 5 & Decay threshold $(s+1)/8$ of the internal peak. \\
Sustain Rate & 0--31; 20 & Reversed chip sustain-decay rate; panel may say RR. \\
\bottomrule
\end{tabularx}

The sustain threshold spans 12.5--100\%, not the 0--100\% suggested by the
current tooltip. It marks the transition into sustain decay, not necessarily
a held plateau. These are discrete register settings, not milliseconds.

GATE and RETRIG detect a rising edge at 2\,V and rearm below 0.01\,V. Either
rising edge retriggers, but GATE's held state decides whether the envelope
remains active. A retrigger without a held gate immediately proceeds toward
release on subsequent samples. Keep the gate high for a full contour.
